\section{Critical regular variation and the divergent branch}
\label{sec:critical}

We now identify the entrywise limit in Theorem~\ref{thm:main}.
The only regular-variation input is the classical uniform convergence
theorem for positive measurable slowly varying functions
\citep{bingham1987regular}: for each compact $J\subset(0,\infty)$,
\begin{equation}
 \sup_{t\in J}\abs{\frac{L(tx)}{L(x)}-1}\longrightarrow0
 \qquad(x\to\infty).
 \label{eq:uct}
\end{equation}
We state the version used explicitly, without requiring monotonicity of
$L$ or convergence of $L$ itself.

\begin{lemma}[Critical partial sums]
\label{lem:critical-sums}
If $F(N)\to\infty$, then
\begin{equation}
 L(N)^2=o(F(N)),\qquad
 \frac{F(N/c)}{F(N)}\longrightarrow1
 \quad\text{for each fixed }c\ge1.
 \label{eq:critical-sums}
\end{equation}
\end{lemma}
\begin{proof}
Fix $A>1$. Uniform convergence on $[1/A,1]$ gives
\begin{equation}
 \frac{1}{L(N)^2}
 \sum_{\lceil N/A\rceil\le k\le N}\frac{L(k)^2}{k}
 \longrightarrow\log A,
 \label{eq:multiplicative-block}
\end{equation}
since the harmonic sum on the same interval tends to $\log A$.
Hence $\liminf_{N\to\infty}F(N)/L(N)^2\ge\log A$.
Only after this limit do we let $A$ grow: for every fixed bound $M$,
choose $A$ with $\log A>M$. It follows that $F(N)/L(N)^2\to\infty$.

For fixed $c>1$, the omitted indices in
$F(N)-F(N/c)$ satisfy $N/c<k\le N$.
By \eqref{eq:uct}, their weights $L(k)^2$ are at most a fixed
multiple of $L(N)^2$ for all sufficiently large $N$, and
$\sum_{N/c<k\le N}1/k$ is bounded in $N$.
Thus $0\le F(N)-F(N/c)=O_c(L(N)^2)=o(F(N))$.
Floors change only the interval endpoints and are already included
in these sums. The case $c=1$ is immediate.
\end{proof}

\begin{lemma}[Finite Gram identity]
\label{lem:gram}
For fixed $m,n$ and $\ell=\lcm(m,n)$, if $N\ge\ell$ then
\begin{equation}
 \ip{e_m}{A_Ne_n}
 =\frac{\sqrt{mn}}{\ell F(N)}
   \sum_{k\le N/\ell}
   \frac{L(k\ell/m)L(k\ell/n)}{k}.
 \label{eq:finite-gram}
\end{equation}
If $N<\ell$, the entry is zero.
\end{lemma}
\begin{proof}
A row contributes to $\ip{T_Ne_m}{T_Ne_n}$ precisely when it is
a common multiple of $m$ and $n$. Write that row as $k\ell$.
Its contribution is
\[
 a(k\ell/m)a(k\ell/n)
 =\frac{\sqrt{mn}}{k\ell}
   L(k\ell/m)L(k\ell/n).
\]
Summing over $k\ell\le N$ and dividing by $F(N)$ proves the claim.
This is the specialization of the arbitrary-coefficient identity in
\citet[Proposition~2.1]{hilberdink2017singular}.
\end{proof}

\begin{proof}[Proof of Theorem~\ref{thm:main}, divergent branch]
Fix $m,n$, and set $\ell=\lcm(m,n)$.
The two factors $\ell/m$ and $\ell/n$ are fixed positive integers, so
slow variation implies
\[
 \frac{L(k\ell/m)L(k\ell/n)}{L(k)^2}\longrightarrow1.
\]
For any $\varepsilon>0$, this ratio lies between $1-\varepsilon$
and $1+\varepsilon$ for all $k$ beyond a fixed index.
Split the sum in \eqref{eq:finite-gram} at that index.
The initial part is finite, whereas $F(N/\ell)\to\infty$.
Consequently
\[
 \sum_{k\le N/\ell}
 \frac{L(k\ell/m)L(k\ell/n)}{k}
 \sim F(N/\ell)\sim F(N),
\]
the second equivalence following from Lemma~\ref{lem:critical-sums}.
Formula~\eqref{eq:finite-gram} therefore tends to
$\sqrt{mn}/\ell=\gcd(m,n)/\sqrt{mn}$.

Unique prime factorization identifies this kernel as
\[
 K(m,n)=\prod_p
 \left(p^{-1/2}\right)^{\abs{v_p(m)-v_p(n)}}.
\]
The classical divergence of $\sum_p1/p$ gives the hypothesis of
Propositions~\ref{prop:singularity} and \ref{prop:collapse}.
They establish pure singularity, the nonexistence of a nonnegative
closed extension, \eqref{eq:divergent-limit}, and norm divergence.
\end{proof}

The diagonal identity
\begin{equation}
 \ip{e_1}{A_Ne_1}
 =F(N)^{-1}\sum_{r\le N}a(r)^2=1
 \label{eq:diagonal-one}
\end{equation}
holds for every $N$, in both branches.
In the divergent branch it precludes strong operator convergence
$A_N\to0$, since such convergence would force the displayed scalar
to tend to zero. Strong-resolvent convergence in
\eqref{eq:divergent-limit} is a different assertion and is compatible
with \eqref{eq:diagonal-one}.
